\scanpage{3}
% L01–02
$\mathbb F_p$, $G_a=G_{a\,\mathbb F_p}$, $n\in\mathbb N$, $X_n$ schéma des sous-groupes finis de rang $p^n$ de $G_a$,
% L02–04
\[
\begin{aligned}
\mathbb E^n=E^n_{\mathbb F_p}&\xrightarrow{\varphi_n}X_n\\
a=(a_0,\ldots,a_{n-1})&\longmapsto\operatorname{Ker}(f_a)\\
f_a:G_a\longrightarrow G_a &\qquad f_a(x)=x^{p^n}+a_{n-1}x^{p^{n-1}}\ \text{---}+a_1x^p+a_0x
\end{aligned}
\]

% L05
\uline{Lemme} $\varphi_n$ est une \deleted{monomorphisme} \inserted{imm. fermée} surjecti\deleted{\unsure{f}}ve, donc induit $\mathbb E^n\xrightarrow{\sim}X_{n\,\unsure{\mathrm{réd}|\mathrm{red}}}$

% L06–07
a) surjectif : il suffit de prouver $k^n\longrightarrow X_n(k)$ surj. si $k$ Alg clos. Or si $\deleted{\unsure{K}}H\subset G_{a\,k}$, $H\in X_n(k)$ i.e. $H$ fini de rang $p^n$, on a
% L08–09
\[
H\simeq\underset{H^0}{\underbrace{\alpha_{p^r}}}\mathbin{\underset{\cdot}{\times}}
\underset{\uparrow\,H\deleted{\unsure{(}}(k)_k}{\left(\mathbb Z/p\right)^s_k}
\qquad r+s=n .
\]
% L08–09
Soit $\xi_1\ldots\xi_s$ base de $H(k)$ sur $\mathbb F_p$, On a $H\deleted{\unsure{Z}}(k)_k=\operatorname{Ker}(\deleted{\unsure{g}}f_a)$, $a\in\mathbb E^s_k$,
% L10
\margin{justifier}
% L10–11
\[
\left\{\begin{aligned}
&\alpha=(\alpha_0,\ldots,\alpha_{s-1})\\
&f_{\unsure{a|\alpha}}(\unsure{x})=x^{p^s}+\alpha_{s-1}x^{p^{s-1}}\ \text{---}+\alpha_1x^p+\alpha_0x
\end{aligned}\right.
\]

% L12
Considérons alors
% L12–14
\[
\begin{aligned}
&f:G_a\mathbin{\text{---}}G_a\qquad f=f_a\operatorname{Frob}^{r}\\
&f(x)=x^{p^{r+s}}+\underset{\substack{\shortparallel\\a_{r+s-1}}}{\alpha_{s-1}}x^{p^{r+s-1}}\ \text{---}
+\underset{\substack{\shortparallel\\a_r}}{\alpha_0}x^{p^r}\\
&\hspace{15em}a_{r-1}=\ldots=a_0=0
\end{aligned}
\]
% L15
\[
\text{On a}\qquad\operatorname{Ker} f=f_a^{-1}\bigl(\underset{\substack{\shortparallel\\\operatorname{Ker}\operatorname{Frob}^r}}{\alpha_{p^r}}\bigr)\supset H(k)_k\mathbin{\cdot}\alpha_{p^r}
\]
% L16–17
et comme les 2 mbs sont de $=$ rang $p^{r+s}=p^n$, l'incl. $\supset$ est égalité.

% L18–20
b) \uline{Mono} Soient \deleted{$a,b\in\unsure{E}$} $k$ $\mathbb F_p$-Alg., $a,b\in\mathbb E^n$ tels que $\varphi_n(a)=\varphi_n(b)$ i.e. $\operatorname{Ker}f_a=\operatorname{Ker}f_b$, prouvons $f_a=f_b$ (i.e. $a=b$) \deleted{\unsure{1\textsuperscript{er} Cas} \unsure{$k$ corps alg. clos}}
% L21–22
\diagram{\begin{tikzcd}[ampersand replacement=\&,column sep=small]
H \arrow[r,phantom,"\subset" description] \& G_{a\,k}\arrow[rr,shift left=2pt,"f_a"]\arrow[rr,shift right=2pt,"f_b"']\arrow[dr] \& \& G_{a\,k}\\
\& \& G_{a\,k}/H\arrow[ur,shift left=2pt,"\unsure{\sim}"]\arrow[ur,shift right=2pt,"\sim"'] \& g_a,g_b
\end{tikzcd}}
% L23–24
$g_b=ug_a$ donc $f_b=uf_a$, $u\in\operatorname{Aut}(G_{a\,k})$. Il faut prouver $u=\mathrm{id}_{G_a}$. Clair si $k$ réduit car $\operatorname{Aut}(G_a)\simeq k^*$.
% archive sheet number: 2
